\documentclass[12pt, a4paper]{article}

% ── Packages ──────────────────────────────────────────────────────────────────
\usepackage[T1]{fontenc}
\usepackage[utf8]{inputenc}
\usepackage{times}
\usepackage[top=1in, bottom=1in, left=1.1in, right=1.1in]{geometry}
\usepackage{fancyhdr}
\usepackage{amsmath}
\usepackage{graphicx}
\usepackage{booktabs}
\usepackage{array}
\usepackage{url}
\usepackage[hidelinks]{hyperref}
\usepackage{setspace}
\usepackage{titlesec}
\usepackage{parskip}
\usepackage{caption}
\usepackage{microtype}
\usepackage{enumitem}
\usepackage{tabularx}
\usepackage{multirow}

% ── Header / Footer ───────────────────────────────────────────────────────────
\pagestyle{fancy}
\fancyhf{}
\lhead{\small\textit{Dr. Karpagam M et al.}}
\rhead{\small\textit{ResearchGPT}}
\renewcommand{\headrulewidth}{0.4pt}
\cfoot{\thepage}
\setlength{\headheight}{14pt}

% ── Section formatting ────────────────────────────────────────────────────────
\titleformat{\section}{\large\bfseries}{\thesection.}{0.5em}{}
\titleformat{\subsection}{\normalsize\bfseries}{\thesubsection}{0.5em}{}
\titlespacing*{\section}{0pt}{2ex plus 0.5ex minus 0.2ex}{1ex}
\titlespacing*{\subsection}{0pt}{1.5ex plus 0.3ex minus 0.1ex}{0.6ex}

% ── Line spacing ──────────────────────────────────────────────────────────────
\setstretch{1.15}

% ── URL font ──────────────────────────────────────────────────────────────────
\urlstyle{tt}

% ══════════════════════════════════════════════════════════════════════════════
\begin{document}

% ── Title block ───────────────────────────────────────────────────────────────
\begin{center}

  {\Huge\bfseries ResearchGPT: An AI-Powered IEEE\\[0.4em]
   Research Paper Generation System\\[0.4em]
   Using Retrieval-Augmented Generation}

  \vspace{1.6em}

  {\large
    Vishal Raghav V$^{1}$,\quad
    Lakshmi Narasimman S$^{2}$,\quad
    Dr.\ Karpagam M$^{3}$}

  \vspace{1em}

  {\small
    $^{1}$Dept.\ of Computing Technologies, SRMIST, Kattankulathur, India\\
    \texttt{vv0636@srmist.edu.in}}

  \vspace{0.5em}

  {\small
    $^{2}$Dept.\ of Computational Intelligence, SRMIST, Kattankulathur, India\\
    \texttt{ls6726@srmist.edu.in}}

  \vspace{0.5em}

  {\small
    $^{3}$Dept.\ of Computational Intelligence, SRMIST, Kattankulathur, India\\
    \texttt{karpagam1@srmist.edu.in}}

\end{center}

\vspace{1em}
\noindent\rule{\linewidth}{0.5pt}
\vspace{0.3em}

% ── Abstract ──────────────────────────────────────────────────────────────────
\begin{center}
  \begin{minipage}{0.88\textwidth}
    \begin{center}{\normalsize\bfseries Abstract}\end{center}
    \vspace{0.4em}
    \small
    The rapid growth of academic literature has created an increasing demand
    for automated tools that can assist researchers in synthesizing and presenting
    knowledge in structured formats. This paper presents ResearchGPT, an end-to-end
    intelligent system that generates IEEE-style research papers from a user-specified
    topic using Retrieval-Augmented Generation (RAG). The system integrates a
    multi-strategy retrieval pipeline combining dense vector search, BM25 sparse
    retrieval, Hypothetical Document Embedding (HyDE), and cross-encoder reranking
    to retrieve high-quality academic context from a curated corpus of arXiv papers.
    A large language model (LLM) hosted on the Cerebras inference platform then
    synthesizes the retrieved context into a structured, multi-section IEEE paper.
    In addition to paper generation, the system provides an automated literature
    survey module that fetches up to fifteen arXiv papers, extracts eight structured
    fields per paper using LLM-based analysis, and presents the results in a
    twelve-column comparative table alongside a generated IEEE-style survey PDF.
    The system further supports post-generation refinement, enabling users to
    iteratively edit specific sections. An evaluation framework based on G-Eval
    assesses faithfulness and relevance of each generated paper. Experimental results
    demonstrate that the HyDE retrieval strategy consistently achieves high relevance
    scores (4--5/5), while faithfulness varies primarily with the quality and
    breadth of available arXiv source coverage. A Streamlit-based user interface
    with multi-session chat management, unified session history, and export options
    provides an accessible front-end for researchers. The complete system, including
    ingestion, retrieval, generation, evaluation, literature survey, and export
    modules, is described in detail.
  \end{minipage}
\end{center}

\vspace{0.8em}
\noindent\textbf{Keywords:} Retrieval-Augmented Generation, Large Language Models,
IEEE Paper Generation, Academic Writing Assistance, Information Retrieval,
Literature Survey Automation, HyDE, Cross-Encoder Reranking
\vspace{0.3em}

\noindent\rule{\linewidth}{0.5pt}

\vspace{1em}

% ══════════════════════════════════════════════════════════════════════════════
\section{Introduction}

The volume of academic publications has grown exponentially over the past decade,
making it increasingly difficult for researchers to stay current with developments
in their respective fields. According to estimates from arXiv.org \cite{arxiv2024},
the repository alone receives over two million new submissions annually, spanning
physics, mathematics, computer science, and related disciplines. Simultaneously,
the task of writing structured research papers demands significant time and expertise
in academic conventions, citation practices, and domain-specific knowledge synthesis.
Automated tools for research assistance have therefore become a focus of considerable
interest in the natural language processing community.

Large language models (LLMs) \cite{brown2020} have demonstrated remarkable
capability in generating coherent, contextually appropriate text across diverse
domains. However, their tendency to hallucinate factual information \cite{ouyang2022}
makes them unreliable when used in isolation for academic writing tasks.
Retrieval-Augmented Generation (RAG) \cite{lewis2020} addresses this limitation
by grounding generation in retrieved documents, reducing hallucination and
improving factual accuracy. Despite the promise of RAG-based approaches, existing
systems have largely focused on question answering or short-form generation rather
than the production of complete, structured academic papers with multiple sections,
inline citations, and evaluation scores.

Beyond paper generation, researchers also require systematic tools for conducting
literature surveys. Traditional literature reviews are labour-intensive processes
that involve searching multiple databases, reading full papers, extracting relevant
information, and synthesizing findings across dozens of sources. Automated extraction
of structured fields from academic papers using LLMs represents a significant
opportunity to accelerate this process, providing researchers with a structured
comparative view of the literature before undertaking original writing.

This paper introduces ResearchGPT, a system that addresses both of these needs
through a unified pipeline. ResearchGPT combines a multi-strategy RAG pipeline
with a state-of-the-art LLM to generate complete, structured IEEE research papers
from a user-provided topic. Uniquely, it also provides an automated literature
survey module that fetches real arXiv papers, extracts structured metadata fields
using LLM-based analysis, and renders the results in a twelve-column table and
a generated survey PDF. Unlike prior work that focuses on summarization or question
answering, ResearchGPT produces fully formatted multi-section papers complete with
abstract, introduction, methodology, results, conclusion, and references.

The primary contributions of this work are: (1) a complete end-to-end pipeline
for IEEE paper generation using RAG; (2) an automated literature survey module
that extracts eight structured fields per paper from live arXiv sources; (3) a
multi-strategy retrieval module supporting standard hybrid retrieval, HyDE
\cite{gao2022}, and complex multi-step retrieval with automatic strategy selection;
(4) a G-Eval-based \cite{liu2023} evaluation framework for assessing faithfulness
and relevance; (5) a post-generation refinement mechanism for iterative editing;
(6) a multi-format export module supporting LaTeX/PDF, DOCX, and BibTeX output;
and (7) a multi-session Streamlit interface supporting concurrent paper generation
workflows with unified session history.

The remainder of this paper is organized as follows. Section~2 reviews related
work. Section~3 describes the overall system architecture. Section~4 details the
literature survey module. Sections~5 through 9 describe the RAG pipeline modules.
Section~10 describes the export and output module. Section~11 presents the user
interface. Section~12 presents experimental results. Section~13 describes
post-generation refinement. Section~14 discusses findings and limitations.
Section~15 concludes.

% ══════════════════════════════════════════════════════════════════════════════
\section{Related Work}

\subsection{Large Language Models for Text Generation}

The emergence of transformer-based language models \cite{devlin2019} has
fundamentally changed how natural language generation tasks are approached.
BERT \cite{devlin2019} introduced bidirectional pre-training, enabling deep
contextual representations that substantially improved performance on a wide range
of downstream tasks. Subsequent models such as GPT-3 \cite{brown2020} demonstrated
that scale alone---in terms of parameters and training data---yields impressive
few-shot generalization without task-specific fine-tuning. InstructGPT
\cite{ouyang2022} further refined generation quality through reinforcement learning
from human feedback (RLHF), producing outputs that are more aligned with user
intent and less prone to toxic or off-topic responses. These models form the
backbone of modern academic writing assistance tools, though their propensity for
hallucination remains a significant challenge for factual domains such as academic
research.

\subsection{Retrieval-Augmented Generation}

RAG \cite{lewis2020} was introduced as a framework for augmenting
sequence-to-sequence models with a non-parametric memory component. By retrieving
relevant documents at inference time and conditioning generation on them, RAG
substantially reduces factual errors compared to purely parametric models and
enables the generation of up-to-date, source-grounded content without retraining.
Since its introduction, many variants have been proposed. Hypothetical Document
Embedding (HyDE) \cite{gao2022} improves dense retrieval by first generating a
hypothetical answer to the query and using its embedding for retrieval, bridging
the semantic gap between short queries and long documents. Self-RAG
\cite{asai2023} introduces adaptive retrieval by training models to decide when
to retrieve and how to use retrieved passages, reducing unnecessary retrieval
overhead. FLARE \cite{jiang2023} performs active retrieval triggered by low
generation confidence, enabling more targeted use of external knowledge.

\subsection{Passage Retrieval and Reranking}

Hybrid retrieval combining sparse BM25 and dense vector search has become a
standard approach in modern information retrieval systems, consistently outperforming
either method in isolation \cite{karpukhin2020}. Cross-encoder
reranking \cite{nogueira2019} refines the initial retrieval set by jointly encoding
query-document pairs, yielding significantly higher precision than bi-encoder
approaches at the cost of additional computation. Sentence-BERT \cite{reimers2019}
provides efficient sentence embeddings suitable for semantic similarity search,
enabling scalable dense retrieval without the latency of full cross-encoder scoring.
MiniLM \cite{wang2020} offers a computationally efficient distilled model that
maintains strong retrieval performance while being practical for deployment on
CPU hardware, which is the target environment for ResearchGPT.

\subsection{Automated Literature Survey and Academic Writing}

Prior systems for automated academic writing have largely focused on abstract
generation, related work synthesis, or paragraph-level generation. Tools such as
Semantic Scholar's AI assistant and Elicit have demonstrated the ability to retrieve
and summarize academic papers on demand, though they do not produce structured
multi-section papers. SurveyAgent and similar tools generate related work sections
by querying paper databases, but do not extend to full paper generation or structured
field extraction. ResearchGPT extends these approaches to full IEEE paper generation
with structured sections, inline citations, evaluation scoring, and a companion
literature survey module that extracts eight structured fields per paper.

\subsection{Evaluation of Natural Language Generation}

Evaluating the quality of LLM-generated text is a well-studied problem. BLEU and
ROUGE provide n-gram overlap metrics suitable for translation and summarization
respectively, but are poorly suited to long-form academic generation where surface
overlap with a reference is not a reliable quality signal. BERTScore \cite{zhang2020} provides a semantic similarity
metric that better captures paraphrase equivalence. G-Eval \cite{liu2023}
represents the current state of the art for reference-free evaluation of long-form
text, using GPT-4 with chain-of-thought prompting to score dimensions such as
coherence, consistency, fluency, and relevance. ResearchGPT adopts a G-Eval-inspired
framework using the Cerebras LLM as the judge, scoring faithfulness and relevance
on a 1--5 scale.

% ══════════════════════════════════════════════════════════════════════════════
\section{System Architecture}

ResearchGPT follows a modular pipeline architecture consisting of seven core
components: (1) Literature Survey, (2) Document Ingestion and Preprocessing,
(3) Embedding and Vector Store, (4) Retrieval, (5) Generation, (6) Evaluation,
and (7) Export. These modules communicate through well-defined interfaces, enabling
independent development and testing of each component. The system exposes two
primary workflows to the user: the \textit{literature survey workflow}, which
fetches and analyses real arXiv papers and presents them in a structured table;
and the \textit{paper generation workflow}, which retrieves relevant context from
the indexed corpus and generates a complete IEEE paper.

\begin{figure}[h]
  \centering
  \includegraphics[width=0.85\linewidth]{figures/architecturediagram.png}
  \caption{High-level architecture of the ResearchGPT system showing the seven
           core modules and their interactions. The literature survey pipeline
           (top) and RAG-based paper generation pipeline (middle and bottom)
           share the arXiv ingestion layer but operate independently.}
  \label{fig:architecture}
\end{figure}

As shown in Figure~\ref{fig:architecture}, a user query enters the system through
the Streamlit UI, triggering two parallel workflows. In the survey workflow, papers
are fetched from arXiv \cite{arxiv2024}, analysed by the LLM to extract eight
structured fields, and displayed in a twelve-column comparative table. In the
generation workflow, the fetched PDFs are chunked and embedded using MiniLM
\cite{wang2020}, and stored in ChromaDB \cite{chroma2024}. At query time, the
retrieval module fetches relevant chunks using the selected strategy, the generation
module produces the structured paper, and the evaluation module scores the output.
Both workflows share the same arXiv ingestion layer, and the results of both are
persisted in the session history for later retrieval.

Table~\ref{tab:modules} summarises the key technologies employed by each module.

\begin{table}[h]
  \centering
  \caption{Module--Technology Mapping for ResearchGPT}
  \label{tab:modules}
  \begin{tabularx}{\linewidth}{lX}
    \toprule
    \textbf{Module} & \textbf{Key Technologies} \\
    \midrule
    Literature Survey   & arXiv API, PyMuPDF, Cerebras LLM \\
    Ingestion           & arXiv API, PyMuPDF, RecursiveCharacterTextSplitter \\
    Embedding           & all-MiniLM-L6-v2, Sentence Transformers, ChromaDB \\
    Retrieval           & BM25 (rank\_bm25), ChromaDB, cross-encoder/ms-marco-MiniLM-L-6-v2 \\
    Generation          & Cerebras API, llama3.1-8b-instruct \\
    Evaluation          & G-Eval (Cerebras LLM judge) \\
    Export              & pdflatex, python-docx, BibtexParser \\
    \bottomrule
  \end{tabularx}
\end{table}

% ══════════════════════════════════════════════════════════════════════════════
\section{Literature Survey Module}

\subsection{Overview}

The literature survey module is a distinct pipeline within ResearchGPT that
provides researchers with a structured comparative analysis of existing work on
their topic before they undertake original writing. Upon receiving a topic, the
module fetches up to fifteen papers from arXiv, downloads and parses their full
text, and uses the LLM to extract eight structured fields from each paper. The
results are presented to the user in a twelve-column interactive table and also
compiled into a generated IEEE-style survey PDF.

\subsection{Paper Fetching and Analysis}

Papers are fetched from arXiv using a keyword-based query derived from the user's
topic, sorted by relevance. The system downloads up to 15 papers by default
(configurable between 3 and 15), with full PDF text extracted via PyMuPDF
\cite{pymupdf2024}. For each paper, the LLM is prompted to extract the following
eight structured fields, with a constraint of one to three sentences per field:

\begin{enumerate}[noitemsep]
  \item \textbf{Problem Statement} --- the concrete problem or challenge addressed
  \item \textbf{Objective} --- the specific goal or contribution
  \item \textbf{Research Gap} --- what was missing in prior work
  \item \textbf{Methodology} --- the exact algorithms, architectures, or techniques
  \item \textbf{Features Focused} --- the specific input features or signals used
  \item \textbf{Results \& Discussion} --- exact metric names and numeric values reported
  \item \textbf{Limitations} --- explicit limitations the authors acknowledge
  \item \textbf{Future Scope} --- concrete future directions suggested
\end{enumerate}

The extraction prompt explicitly forbids vague paraphrasing, requiring the LLM to
quote or closely paraphrase actual content from the paper. Fields not present in
the paper are recorded as ``Not mentioned'' rather than hallucinated. Up to 7000
characters of extracted PDF text are passed to the LLM per paper to keep inference
time manageable.

\subsection{Structured Survey Table}

The extracted fields, combined with bibliographic metadata (title, authors, year,
DOI/arXiv ID, citation string), form a twelve-column structured table. The table
supports two display modes in the UI: a \textit{Table view} with sticky headers
and columns for scrollable comparison, and a \textit{Card view} displaying each
paper as an expandable card with a three-column grid showing Problem Statement,
Objective, and Research Gap. An arXiv link column provides direct access to the
source paper for each entry.

\subsection{Survey PDF Generation}

In addition to the interactive table, the module generates a formal IEEE-style
survey paper summarising the analysed papers. The survey paper is produced by
prompting the LLM with a condensed context derived from the extracted fields of
the top ten papers, and instructed to produce a complete multi-section survey
including Abstract, Introduction, Literature Review, Problem Definition,
Methodology, Results, Discussion, and Conclusion, with inline citations and a
References section in IEEE format. This survey paper is available for download
directly from the UI.

% ══════════════════════════════════════════════════════════════════════════════
\section{Document Ingestion and Preprocessing}

\subsection{Paper Acquisition for RAG}

For the RAG-based paper generation pipeline, papers are acquired dynamically from
the arXiv API \cite{arxiv2024} based on the user's research topic. The system
submits a relevance-sorted keyword-based search query and downloads up to five PDF
papers per topic for indexing. A rate limiter enforces a one-second delay between
consecutive downloads to comply with arXiv's usage policy. Failed downloads are
retried up to three times with exponential backoff. If a PDF download fails
entirely, the paper's abstract is used as a fallback text source to ensure partial
coverage even when full text is unavailable.

\subsection{PDF Parsing}

Downloaded PDFs are parsed using PyMuPDF \cite{pymupdf2024}, which extracts
raw text while preserving paragraph structure. Post-extraction, a cleaning pass
removes headers, footers, page numbers, and \LaTeX{} artefacts using regular
expression filters. Text is normalized by collapsing whitespace, removing control
characters, and standardizing quotation marks. Up to ten pages are extracted per
paper to balance content coverage against memory usage.

\subsection{Chunking}

Cleaned document text is split into overlapping chunks of 1000 characters with a
350-character overlap. Overlap ensures that sentences spanning chunk boundaries
are represented in both adjacent chunks, reducing information loss at boundaries.
Each chunk is stored alongside metadata including the source paper title, authors,
and arXiv identifier, enabling provenance tracking during retrieval. The chunk
metadata is subsequently used to build the references section of the generated
paper, ensuring that all cited sources correspond to real retrieved documents.

% ══════════════════════════════════════════════════════════════════════════════
\section{Embedding and Vector Store}

Chunk embeddings are computed using the \texttt{all-MiniLM-L6-v2} model
\cite{wang2020} from the Sentence Transformers library \cite{reimers2019}. This
model produces 384-dimensional dense vectors and offers a favorable balance between
embedding quality and inference speed on CPU hardware, making it well-suited for
deployment in resource-constrained environments. The model achieves strong
performance on semantic textual similarity benchmarks while requiring no GPU
hardware for inference.

Embeddings are stored and indexed in ChromaDB \cite{chroma2024}, a lightweight
open-source vector database designed for AI applications. ChromaDB provides a
persistent collection interface that supports cosine similarity search with
efficient approximate nearest-neighbour indexing. Collections are named by topic
slug (a normalized, lowercase, alphanumeric identifier derived from the topic
string) and are reused across sessions for the same topic to avoid redundant
re-indexing. A stale collection eviction policy ensures that when the number of
stored collections exceeds ten, the oldest collection is deleted to reclaim disk
storage and prevent unbounded growth over long-running deployments.

BM25 index construction is performed in parallel with embedding. The BM25 index
is built over the tokenized chunk texts using the \texttt{rank\_bm25} library,
providing a complementary sparse retrieval index that excels at exact keyword
matching. The BM25 and dense indices together form the backbone of the hybrid
retrieval system described in Section~7.

% ══════════════════════════════════════════════════════════════════════════════
\section{Retrieval Strategies}

ResearchGPT exposes four retrieval strategies. Each strategy retrieves a ranked
list of the top-$k$ chunks (configurable, default $k = 6$) for use in generation.
The strategy can be selected by the user or determined automatically by the system.

\subsection{Standard Hybrid Retrieval}

The standard strategy combines BM25 sparse retrieval with dense vector similarity
search. Candidate chunks are retrieved independently from both indices, and their
scores are interpolated with weight $\alpha$ (Vector 70\%, BM25 30\% by default):

\begin{equation}
  \text{score}(d) = \alpha \cdot \text{sim}_{\text{dense}}(d)
                  + (1 - \alpha) \cdot \hat{s}_{\text{BM25}}(d)
  \label{eq:hybrid}
\end{equation}

where $\hat{s}_{\text{BM25}}$ denotes the min-max normalised BM25 score. The
merged candidate set (typically $3k$ candidates) is then passed through a
cross-encoder reranker \cite{nogueira2019} using the
\texttt{cross-encoder/ms-marco-MiniLM-L-6-v2} model, which jointly encodes each
query-document pair to produce a relevance score. The top-$k$ chunks by reranker
score form the final retrieval result. This strategy is the most computationally
efficient and performs well for focused, well-defined queries.

\subsection{HyDE Retrieval}

HyDE \cite{gao2022} addresses the vocabulary mismatch problem between short user
queries and long academic passages. A hypothetical document is first generated by
the LLM for the query topic, producing a passage that resembles the kind of text
that would appear in a relevant paper. The embedding of this hypothetical document
is then used to query the vector store instead of the raw query embedding. This
shifts the retrieval anchor from a sparse query representation to a dense
passage-like representation, substantially improving recall for complex academic
topics where the query itself contains few of the technical terms present in
relevant passages. The retrieved candidates are subsequently reranked by the
cross-encoder using the original query to ensure precision.

\subsection{Complex Multi-Step Retrieval}

The complex strategy decomposes the user query into multiple sub-questions using
the LLM. Typically three to five sub-questions are generated, each targeting a
distinct facet of the original query (e.g., motivations, methods, results,
limitations, future directions). Each sub-question is independently retrieved
against the vector store using dense similarity search. Results from all
sub-queries are merged, deduplicated by chunk ID, and reranked jointly. This
approach is particularly effective for broad or multi-faceted queries where a
single retrieval pass may miss relevant perspectives. Its primary drawback is
latency, as multiple LLM calls and retrieval passes are required.

\subsection{Auto Strategy Selection}

The auto strategy automatically selects between HyDE and complex retrieval based
on query characteristics. Queries containing comparative keywords (e.g.,
\textit{compare}, \textit{versus}, \textit{difference between}, \textit{benchmark})
or exceeding sixteen words in length are routed to the complex strategy; all others
are routed to HyDE. This provides a sensible default for users who do not wish to
manually select a strategy, while still benefiting from the complementary strengths
of both approaches.

% ══════════════════════════════════════════════════════════════════════════════
\section{Generation Module}

\subsection{LLM Backend}

The generation module uses the Cerebras inference platform with the
\texttt{llama3.1-8b-instruct} model via an OpenAI-compatible REST API
\cite{brown2020}. The Cerebras wafer-scale engine delivers significantly higher
token throughput than GPU-based inference --- typically exceeding 1000 tokens per
second --- enabling the generation of long-form academic papers with low end-to-end
latency. Temperature is set to 0.2 to favour factual, consistent outputs over
creative variation. Table~\ref{tab:config} summarises the generation configuration.

\begin{table}[h]
  \centering
  \caption{LLM Generation Configuration}
  \label{tab:config}
  \begin{tabular}{ll}
    \toprule
    \textbf{Parameter} & \textbf{Value} \\
    \midrule
    Provider            & Cerebras              \\
    Model               & llama3.1-8b-instruct  \\
    Temperature         & 0.2                   \\
    Max tokens          & 8192                  \\
    Context chars/chunk & 3000                  \\
    Max chunks (full)   & 9                     \\
    Max chunks (fast)   & 6                     \\
    Target paper length & 13--15 pages          \\
    \bottomrule
  \end{tabular}
\end{table}

\subsection{Prompt Design}

The generation prompt instructs the model to act as a university researcher writing
an IEEE conference paper, outputting plain academic text with Roman numeral section
headings. The prompt is structured in three parts: (1) a system role definition
specifying the academic writing persona and output constraints; (2) the retrieved
context, consisting of up to nine chunk passages each prefixed with their source
paper title and arXiv ID; and (3) a structured generation instruction specifying
section headings, length guidance (150--600 words per section), citation style,
and forbidden elements.

Forbidden elements in the prompt include: markdown formatting, bullet points or
numbered lists within section bodies, fabricated statistics, citations not present
in the provided context, and section headings not listed in the instruction.
The prompt explicitly requests that every factual claim be followed by an inline
citation of the form \texttt{[N]} referring to one of the retrieved sources. A
minimum page target of thirteen pages is enforced through post-generation length
checks, with the model re-prompted to expand if the output falls below the minimum.

\subsection{Post-Processing}

After generation, a markdown stripping pass removes any bold, italic, heading, or
list markers that the model may have inserted despite instructions. A citation
validation pass removes any numeric citation markers whose index exceeds the number
of retrieved sources. A humanisation pass applies light paraphrasing to paragraphs
detected as highly similar to retrieved passages, reducing potential plagiarism
scores. Finally, a ROUGE-L similarity check is performed between the generated text
and the retrieved context to flag and log any sections with unusually high overlap.

% ══════════════════════════════════════════════════════════════════════════════
\section{Evaluation Module}

\subsection{G-Eval Framework}

ResearchGPT evaluates each generated paper using a G-Eval \cite{liu2023} inspired
framework that scores two dimensions: \textit{faithfulness} (whether all claims in
the paper are supported by the retrieved context) and \textit{relevance} (whether
the paper substantively addresses the user's topic). Both dimensions are scored on a
1--5 integer scale by the Cerebras LLM acting as a judge. The judge receives the
generated paper text (up to 7000 characters), up to three retrieved context chunks,
and the user's original topic, and is prompted to reason step by step before
assigning a score.

\subsection{Faithfulness Scoring}

The faithfulness prompt instructs the judge to identify any claims in the generated
paper that cannot be verified from the provided context. A score of 5 indicates
that all claims are fully grounded; a score of 1 indicates extensive hallucination.
The judge is also prompted to provide a one-sentence reasoning string explaining
its score, which is displayed in the UI alongside the numeric score.

\subsection{Relevance Scoring}

The relevance prompt instructs the judge to assess whether the generated paper's
content is directly and substantively relevant to the user's stated topic. A score
of 5 indicates that all sections are on-topic; a score of 1 indicates that the paper
largely addresses a different topic. Relevance scores are typically more stable than
faithfulness scores across topics, as the retrieval pipeline's topic-scoped indexing
ensures that retrieved chunks are generally on-topic.

% ══════════════════════════════════════════════════════════════════════════════
\section{Export and Output Module}

ResearchGPT supports three export formats: LaTeX/PDF, DOCX, and plain text. These
are accessible from a persistent download bar at the bottom of the paper view, with
a separate download option for the survey PDF in the literature survey view.

\subsection{LaTeX and PDF Export}

The generated paper text is formatted into a LaTeX document using a custom template
that replicates IEEE conference paper style, including two-column layout, IEEE
fonts, and a formatted references section. BibTeX entries for all cited sources are
constructed from the retrieved chunk metadata (title, authors, year, arXiv ID) and
deduplicated by a normalized citation key. The LaTeX source is compiled using
\texttt{pdflatex} on the server, and the resulting PDF is made available for
download. Compilation errors are caught and logged, with the raw LaTeX source also
available as a fallback.

\subsection{DOCX Export}

A Word-compatible DOCX file is generated using the \texttt{python-docx} library.
The DOCX export preserves the section structure of the generated paper, with
section headings rendered as Heading~1 style and body text rendered as Normal style.
Inline citations are retained in the text body. A References section is appended
with numbered entries corresponding to the cited sources. The DOCX format is
particularly useful for researchers who wish to make further edits in a word
processor before submission.

\subsection{BibTeX Export}

A standalone BibTeX file is generated for all sources retrieved during the paper
generation session. Citation keys are constructed from the first author's last name
and the publication year, with a numeric suffix appended if a key collision occurs.
This file can be imported directly into reference managers such as Zotero or used
with Overleaf for LaTeX compilation.

% ══════════════════════════════════════════════════════════════════════════════
\section{User Interface}

ResearchGPT provides a Streamlit-based web interface with a multi-session chat
architecture modelled on the interaction patterns of modern conversational AI tools
such as Claude and ChatGPT. The interface is designed to minimise friction for
researchers unfamiliar with RAG systems while exposing sufficient control for
advanced users.

\begin{figure}[h]
  \centering
  \includegraphics[width=0.88\linewidth]{figures/Homepg ResearchGPT.png}
  \caption{ResearchGPT home screen showing the topic input area, example topic
           pills, and the Generation Settings button in the sidebar.}
  \label{fig:welcome}
\end{figure}

\begin{figure}[h]
  \centering
  \includegraphics[width=0.88\linewidth]{figures/Output.png}
  \caption{ResearchGPT Quality \& Sources view showing G-Eval faithfulness and
           relevance scores, pipeline timing breakdown, and retrieval metrics
           for the topic \textit{Network Intrusion Detection}.}
  \label{fig:paperview}
\end{figure}

\subsection{Home Screen and Topic Input}

On first load, the interface presents a welcome message from the ResearchGPT
assistant explaining the system's capabilities, followed by three clickable example
topic pills (\textit{finetuning methods for LLMs}, \textit{quantum computing},
\textit{transformer architectures in NLP}). Clicking a pill pre-fills the topic
input and initiates generation. The topic input is implemented as a persistent
Streamlit chat input field anchored to the bottom of the screen, matching the
interaction pattern of modern chat interfaces.

\subsection{Sidebar and Session Management}

The sidebar displays a ``New Research'' button that creates a new named session,
followed by a unified ``Recent'' panel listing all previous sessions with their
topic names. Active sessions are highlighted with a distinctive colour; completed
sessions display a document icon. Each session entry includes a delete button for
removing unwanted history entries. A ``Generation Settings'' button at the bottom
of the sidebar opens a dedicated settings page.

\subsection{Generation Settings}

The settings page, accessible from the sidebar, allows the user to configure four
generation parameters: (1) retrieval strategy (auto / standard / HyDE / complex);
(2) number of survey papers (3--15); (3) top-$k$ retrieval chunks (3--12); and
(4) whether to run G-Eval scoring after generation. Settings are persisted in the
Streamlit session state across page interactions and saved when the user clicks the
``Save \& return'' button.

\subsection{Paper View and Tabs}

Once a paper is generated, the main content area displays three tabs:
\textit{Paper}, \textit{Survey \& Analysis}, and \textit{Quality \& Sources}.
The Paper tab renders the generated IEEE text. The Survey \& Analysis tab presents
the twelve-column literature survey table with a Table/Card view toggle and a
per-paper arXiv link. The Quality \& Sources tab displays faithfulness and relevance
score cards with colour-coded severity indicators (green for good, yellow for fair,
red for poor), the pipeline timing breakdown, and retrieval scores including chunks
retrieved, papers used, and average, maximum, and minimum retrieval scores.

A faithfulness alert banner is shown at the top of the Quality tab when the
faithfulness score falls below 3, informing the user that the generated paper may
contain claims not fully grounded in the retrieved context. A warning badge
(\textbf{!}) is also appended to the Quality \& Sources tab label itself when
faithfulness falls below 3, making the issue visible without requiring the user to
first navigate to that tab. A small colour-coded faithfulness badge (Good / Fair /
Poor) is rendered inline in the paper metadata caption so users can see quality at
a glance from the Paper tab. A section jump navigation bar with clean Roman-numeral
labels (e.g., \textit{I. Introduction}, \textit{II. Methodology}) provides quick
in-paper navigation for long generated papers. The entire interface uses a cohesive
blue colour scheme (\texttt{\#2563eb} primary) replacing the default green theme.
A download bar at the bottom of the Paper tab provides one-click export to PDF,
DOCX, LaTeX source, and plain text.

% ══════════════════════════════════════════════════════════════════════════════
\section{Experiments and Results}

\subsection{Evaluation Methodology}

Generated papers are evaluated using the G-Eval framework described in Section~8.
Each topic is processed end-to-end: arXiv papers are fetched and indexed, a paper
is generated using the specified strategy, and the output is scored by the LLM
judge. Timing measurements are captured for each pipeline stage (indexing,
retrieval, generation) and summed to give total wall-clock time.

\subsection{Retrieval Strategy Comparison}

Table~\ref{tab:strategies} presents evaluation scores for three retrieval strategies
applied to the topic \textit{Network Intrusion Detection}, a well-represented
topic on arXiv with strong PDF availability.

\begin{table}[h]
  \centering
  \caption{Retrieval Strategy Evaluation (Topic: Network Intrusion Detection)}
  \label{tab:strategies}
  \begin{tabular}{lcccc}
    \toprule
    \textbf{Strategy} & \textbf{Faithfulness} & \textbf{Relevance}
                      & \textbf{Retrieval (s)} & \textbf{Total (s)} \\
    \midrule
    Standard & 2/5 & 5/5 & 18.72 & 40.55  \\
    Complex  & 3/5 & 5/5 & 214.38 & 259.64 \\
    HyDE     & 4/5 & 5/5 & 21.06 & 279.48 \\
    \bottomrule
  \end{tabular}
\end{table}

HyDE achieves the highest faithfulness score (4/5) with retrieval time comparable
to the standard strategy (21 seconds), making it the recommended default for
most queries. The complex strategy improves faithfulness over standard but incurs
substantially higher latency due to multiple LLM sub-query calls. The auto strategy
selects HyDE for this topic (no comparative keywords, query length $<$ 16 words)
and thus produces identical scores to the HyDE row.

\subsection{Per-Query Results Across Topics}

Table~\ref{tab:perquery} presents per-query evaluation results across five topics
tested during system development, all using the HyDE strategy.

\begin{table}[h]
  \centering
  \caption{Per-Query Evaluation Results (Strategy: HyDE)}
  \label{tab:perquery}
  \begin{tabular}{lccccc}
    \toprule
    \textbf{Topic} & \textbf{Faith.} & \textbf{Rel.}
                   & \textbf{Papers} & \textbf{Chunks} & \textbf{Total (s)} \\
    \midrule
    Network Intrusion Detection      & 4/5 & 5/5 & 2/5 & 6/6 & 279.48  \\
    Transformer Architectures in NLP & 5/5 & 5/5 & 5/5 & 6/6 & 179.51  \\
    Federated Learning Privacy       & 1/5 & 5/5 & 1/5 & 6/6 & 364.37  \\
    Finetuning Methods for LLMs      & 1/5 & 4/5 & 1/5 & 6/6 & 1135.07 \\
    Graph Neural Networks            & 3/5 & 4/5 & 3/5 & 6/6 & 881.19  \\
    \bottomrule
  \end{tabular}
\end{table}

\subsection{Analysis of Results}

Several patterns are apparent from the results. First, \textit{relevance} remains
consistently high (4--5/5) across all topics, confirming that the HyDE strategy
reliably produces on-topic papers regardless of the specific subject area. The
topic-scoped ChromaDB collection and HyDE's passage-like query embedding together
ensure that retrieved chunks are topically aligned with the user's query.

Second, \textit{faithfulness} correlates strongly with the number of papers
contributing chunks to generation (the ``Papers'' column, reported as papers used
out of papers indexed). Topics where chunks from multiple papers are retrieved
(e.g., \textit{Transformer Architectures in NLP} at 5/5 papers, faithfulness 5/5)
achieve near-perfect faithfulness, as the LLM has a rich, diverse retrieved context
to draw from. Topics where only one paper contributes usable chunks
(e.g., \textit{Federated Learning Privacy} and \textit{Finetuning Methods for
LLMs} at 1/5 papers) achieve low faithfulness scores (1/5), as the model is
forced to rely on its parametric knowledge to fill section content, leading to
claims not grounded in the retrieved context.

Third, \textit{total generation time} varies widely (179--1135 seconds), primarily
driven by indexing time, which is proportional to the number and size of PDFs
successfully downloaded. The retrieval and generation stages contribute modestly to
total time (typically 20--40 seconds combined). These findings indicate that the
primary bottleneck for both faithfulness quality and latency is arXiv PDF
availability and download success rate.

\subsection{Retrieval Score Analysis}

For all tested topics, the retrieval precision (Average/Max/Min reranker scores)
is uniformly 1.000/1.000/1.000, indicating that the cross-encoder consistently
assigns maximum confidence to the top-$k$ retrieved chunks. All six configured
retrieval slots (Chunks: 6/6) are filled in every run, confirming that the indexed
corpus provides sufficient coverage for the retrieval module to operate at full
capacity even when only one or two PDFs are successfully downloaded.

% ══════════════════════════════════════════════════════════════════════════════
\section{Post-Generation Refinement}

ResearchGPT supports iterative post-generation refinement, enabling users to
modify specific aspects of a generated paper through natural language instructions.
After a paper is generated, the user can type an instruction such as \textit{expand
the methodology section} or \textit{add more detail about LoRA fine-tuning} into
the chat input. The refinement module constructs a targeted prompt containing the
original paper, the user instruction, and a sample of the retrieved context, then
calls the LLM to produce a revised version.

The refined paper replaces the previous version in the session state, and the
G-Eval faithfulness and relevance scores are recomputed to reflect the updated
content. The chat history is preserved, allowing users to trace the evolution of
the paper across multiple refinement passes. An iterative pass counter is
maintained in session state to prevent unbounded refinement loops in automated
testing.

This workflow mirrors the iterative drafting process used by human authors and
provides a practical mechanism for improving paper quality without regenerating
from scratch. In practice, refinement instructions targeting specific sections
(e.g., expanding a single section) tend to produce higher-quality results than
broad instructions, as the LLM can focus its attention on a delimited portion of
the paper. Users seeking to substantially restructure the paper are advised to
initiate a new generation with a modified topic or strategy rather than relying
solely on refinement.

% ══════════════════════════════════════════════════════════════════════════════
\section{Discussion}

\subsection{Effectiveness of Retrieval Strategies}

The experimental results demonstrate that HyDE consistently outperforms standard
hybrid retrieval on faithfulness while maintaining comparable latency, making it
the recommended default strategy for most academic topics. The performance advantage
of HyDE stems from its use of a hypothetical document embedding as the retrieval
anchor, which is semantically closer to the dense representations of academic
passages than a short query embedding. For queries containing comparative or
multi-faceted language, the complex strategy offers further benefits through
sub-question decomposition, at the cost of additional latency.

\subsection{Faithfulness and Source Coverage}

The primary determinant of faithfulness quality in ResearchGPT is the breadth of
arXiv source coverage for a given topic. When full PDF text is available for
multiple papers, the retrieved context is sufficiently rich and diverse that the
LLM can write almost entirely from retrieved evidence, achieving near-perfect
faithfulness. When PDF downloads fail and only one paper's text is indexed, the LLM
must rely on parametric knowledge for a substantial portion of the content, leading
to low faithfulness scores even with high relevance. This dependency on PDF
availability is the most significant limitation of the current system.

Several mitigations are possible. First, falling back to abstract text when PDF
download fails (already implemented) provides partial coverage. Second, integrating
additional academic databases such as Semantic Scholar and PubMed would increase
the probability of successful full-text retrieval for any given topic. Third,
increasing the number of arXiv papers fetched (up to the configured maximum of
fifteen) would improve coverage at the cost of higher indexing time.

\subsection{Scalability and Deployment}

The current implementation is designed for single-user local deployment on standard
laptop hardware. The use of ChromaDB as a persistent vector store and the stale
collection eviction policy ensure that storage usage remains bounded over time.
For multi-user deployment, the session state management would need to be adapted
to a server-side persistence model (e.g., Redis or a relational database), and
the synchronous arXiv fetch-and-index pipeline would benefit from asynchronous
execution to avoid blocking concurrent users.

\subsection{Ethical Considerations}

ResearchGPT is designed as an \textit{assistance} tool, not a replacement for
human authorship. The generated papers are explicitly labelled as AI-generated
outputs and are intended to serve as drafts or starting points for human
researchers rather than as final submissions. The G-Eval faithfulness score
provides users with an indicator of how much the generated content is grounded
in retrieved evidence versus potentially fabricated, enabling informed decision
making about which sections require the most careful human review.

% ══════════════════════════════════════════════════════════════════════════════
\section{Conclusion}

This paper presented ResearchGPT, a complete end-to-end system for generating
IEEE-style research papers using retrieval-augmented generation. The system
integrates seven functional modules: literature survey, document ingestion,
embedding and vector storage, multi-strategy retrieval (standard hybrid, HyDE,
complex, and auto), LLM-based generation via Cerebras, G-Eval evaluation,
and multi-format export. An automated literature survey module complements the
paper generation pipeline by providing structured eight-field extraction and a
twelve-column comparative table for up to fifteen arXiv papers per topic.

Experimental evaluation across five diverse topics demonstrates that relevance
is consistently high (4--5/5) across all retrieval strategies, while faithfulness
is primarily determined by the breadth of arXiv PDF coverage for the topic.
The HyDE strategy achieves the best faithfulness-to-latency trade-off among the
evaluated strategies. The G-Eval evaluation framework and per-session quality
indicators provide users with actionable feedback for iterative improvement.

The system's modular architecture supports straightforward extension to new
retrieval strategies, LLM backends, and document corpora. Future work will focus
on integrating additional academic databases, supporting multi-document indexing
across sessions, and investigating larger LLM backends for improved faithfulness
on low-coverage topics. ResearchGPT represents a practical and accessible tool
for researchers seeking AI-assisted paper generation grounded in retrieved academic
evidence.

% ══════════════════════════════════════════════════════════════════════════════
\begin{thebibliography}{99}

\bibitem{devlin2019}
J. Devlin, M.-W. Chang, K. Lee, and K. Toutanova,
``BERT: Pre-training of Deep Bidirectional Transformers for Language Understanding,''
in \textit{Proc. NAACL-HLT}, 2019, pp. 4171--4186.

\bibitem{brown2020}
T. B. Brown \textit{et al.},
``Language Models are Few-Shot Learners,''
in \textit{Advances in Neural Information Processing Systems}, vol. 33, 2020,
pp. 1877--1901.

\bibitem{ouyang2022}
L. Ouyang \textit{et al.},
``Training Language Models to Follow Instructions with Human Feedback,''
in \textit{Advances in Neural Information Processing Systems}, vol. 35, 2022,
pp. 27730--27744.

\bibitem{lewis2020}
P. Lewis \textit{et al.},
``Retrieval-Augmented Generation for Knowledge-Intensive NLP Tasks,''
in \textit{Advances in Neural Information Processing Systems}, vol. 33, 2020,
pp. 9459--9474.

\bibitem{gao2022}
L. Gao, X. Ma, J. Lin, and J. Callan,
``Precise Zero-Shot Dense Retrieval without Relevance Labels,''
\textit{arXiv preprint arXiv:2212.10496}, 2022.

\bibitem{nogueira2019}
R. Nogueira and K. Cho,
``Passage Re-ranking with BERT,''
\textit{arXiv preprint arXiv:1901.04085}, 2019.

\bibitem{liu2023}
Y. Liu \textit{et al.},
``G-Eval: NLG Evaluation using GPT-4 with Better Human Alignment,''
in \textit{Proc. EMNLP}, 2023, pp. 2511--2522.

\bibitem{arxiv2024}
arXiv.org,
``arXiv API Access,''
\url{https://info.arxiv.org/help/api/index.html}, 2024.

\bibitem{pymupdf2024}
Artifex Software,
``PyMuPDF: Python Bindings for MuPDF,''
\url{https://pymupdf.readthedocs.io}, 2024.

\bibitem{chroma2024}
Chroma Core,
``Chroma: the AI-native open-source embedding database,''
\url{https://www.trychroma.com}, 2024.

\bibitem{wang2020}
W. Wang \textit{et al.},
``MiniLM: Deep Self-Attention Distillation for Task-Agnostic Compression
of Pre-Trained Transformers,''
in \textit{Advances in Neural Information Processing Systems}, vol. 33, 2020,
pp. 5776--5788.

\bibitem{reimers2019}
N. Reimers and I. Gurevych,
``Sentence-BERT: Sentence Embeddings using Siamese BERT-Networks,''
in \textit{Proc. EMNLP-IJCNLP}, 2019, pp. 3982--3992.

\bibitem{asai2023}
A. Asai \textit{et al.},
``Self-RAG: Learning to Retrieve, Generate, and Critique through Self-Reflection,''
\textit{arXiv preprint arXiv:2310.11511}, 2023.

\bibitem{jiang2023}
Z. Jiang \textit{et al.},
``Active Retrieval Augmented Generation,''
in \textit{Proc. EMNLP}, 2023, pp. 7969--7992.

\bibitem{karpukhin2020}
V. Karpukhin \textit{et al.},
``Dense Passage Retrieval for Open-Domain Question Answering,''
in \textit{Proc. EMNLP}, 2020, pp. 6769--6781.

\bibitem{zhang2020}
T. Zhang \textit{et al.},
``BERTScore: Evaluating Text Generation with BERT,''
in \textit{Proc. ICLR}, 2020.

\end{thebibliography}

\end{document}
